\section{Simulation study}\label{sec:simulation-study}

To illustrate the power benefit and small sample properties of our proposed procedure we consider four different scenario through a simulation study. In all scenarios it is assumed that treatment (denoted $R_i$) is randomized one to one between two groups, and we let the observed combined outcome be given by the product $\widetilde{Y}_i = A_i Y_i$, where $A_i$ is binary and $Y_i$ is continuously distributed. This is equivalent to observing the atom $\mathcal{E} = 0$ when $A_i = 0$ and otherwise $Y_i$ is observed.

Table~\ref{tab:simulation} shows the distributions of $A_i \mid R_i$ and $Y_i \mid A_i = 1, R_i$ for each of the four simulation scenarios. Scenario 1 represents the case where the point mass at zero is independent of the treatment but with a continuous component that is shifted between the groups. Scenario 2 is the opposite of scenario 1 where the continuous component is independent of the treatment but the probability of the point mass at zero differs between treatment groups. In scenario 3 we model a treatment effect on both components but in opposite directions of each other, and scenario 4 is similar to scenario 1 but where the continuous component is non-normal distributed.

\begin{table}[htbp]
\centering
\begin{tabular}{r|r|r}
Scenario & $A_i \mid R_i$ & $Y_i \mid A_i = 1, R_i$\\ \hline
1 & $\Bernoulli(0.35)$ & $N(3 + 0.5R_i, 1)$\\
2 & $\Bernoulli(0.5 + 0.15R_i)$ & $N(3.5, 1)$\\
3 & $\Bernoulli(0.4 - 0.1R_i)$ & $N(3 + 0.5 R_i, 1)$\\
4 & $\Bernoulli(0.35)$ & $\Beta(1, 1-0.7R_i)$
\end{tabular}
\caption{Power simulation scenarios}
\label{tab:simulation}
\end{table}

To assess the power we vary the sample size between 50 and 350 observations in each group with increments of 25, and for each scenario we perform 25,000 Monte Carlo simulations and estimate the power as the proportion of rejected null hypotheses at the 5\% level. We compare the estimated power for both our parametric and semi-parametric likelihood ratio tests to the two-sample t-test and the Wilcoxon test. Figure~\ref{fig:powerCurves} shows the estimated power functions. The results are also presented in table form in Appendix~\ref{app:supplementary-simulation-results}.

In all scenarios except number two it is clear that we observe a large power gain compared to both the t-test and the Wilcoxon test.

In setup 1 we observe a large power gain compared to the Wilcoxon test. Here the Wilcoxon tests gets ``confused'' by the large number of ties in the atom. In setup 2 Wilcoxon test has slightly better power profile. This is to be expected as our novel test is here disadvantaged by being a two-degrees-of-freedom test where the Wilcoxon is only one. It is further observed that Wilcoxon as expected as very little power when the effects on mortality and among survivors are of opposite sign despite the two distributions being markedly different indicating clear treatment effect (setups 3 and 4). In contrast our novel methods has excellent power. In all settings with a normally distributed outcome among survivors the parametric and semi-parametric approaches are similar, but for the heavy tail setup (no.~4) the semi-parametric approach is clearly superior.

\begin{figure}[htbp]
\centering
\includegraphics[width=0.84\textwidth]{power-curves.pdf}
\caption{Simulated power as a function of sample size for each of the four scenarios.}
\label{fig:powerCurves}
\end{figure}

In addition to the power simulations we also performed a simulation study for the type I error. We considered two scenarios similar to scenarios 1 and 4 in Table~\ref{tab:simulation} under a null hypothesis of no treatment effect. Specifically, we considered $A_i \sim \Bernoulli(\pi)$, $Y_i \mid A_i = 1 \sim N(3,1)$ and $A_i \sim \Bernoulli(\pi)$, $Y_i \mid A_i = 1 \sim \Beta(1, 1)$. We estimated the type I error through Monte Carlo simulation for different sample sizes and $\pi \in \left\{0.2, 0.4, 0.6, 0.8\right\}$. The results are presented in Appendix~\ref{app:supplementary-simulation-results}.
